% --- Archon physics formalization source begin ---
% archon:physics
% archon:covers PhyXMiniProblems/problem_phyx_mini_0070.lean
% archon:source-report reports/phyx_mini/problem_phyx_mini_0070.source.json

\chapter{Physics problem phyx\_mini\_0070}
\label{ch:PhyXMiniProblems_problem_phyx_mini_0070}

\paragraph{Problem source.}
\#\# Physical scenario

The meter stick in Figure lies on the bottom of a 100-cm-long tank with its zero mark against the left edge. You look into the tank at a 30\textasciicircum{}\textbackslash{}circ angle, with your line of sight just grazing the upper left edge of the tank.

\#\# Question

What mark do you see on the meter stick if the tank is empty?

\#\# Answer choices

- A: A: 24cm

- B: B: 56cm

- C: C: 87cm

- D: D: 99cm

\#\# Figure caption (auxiliary; use the image as primary evidence)

The image depicts a setup involving a meter stick positioned horizontally along the bottom of a rectangular container. Here are the details:

- **Objects and Measurements:**
  - The meter stick is labeled at its starting point with "Zero."
  - The horizontal length of the setup, where the meter stick is placed, is labeled as "100 cm."
  - The vertical height of the container is labeled as "50 cm."

- **Line of Sight:**
  - An angled line labeled "Line of sight" forms a 30-degree angle with the vertical side of the container. This suggests viewing or measuring from that angle.

- **Layout:**
  - The setup includes a rectangular container that appears to hold the meter stick, emphasizing distance and angles for measurement or viewing purposes.

\#\# Dataset metadata

Geometrical Optics; Spatial Relation Reasoning

\paragraph{Recorded answer/context.}
C: C: 87cm

\paragraph{Figure/image path.}
/root/proposal\_for\_physic/science-mango/phyx\_mini\_run/phyx\_data/test\_image/70.png

\paragraph{Formalization target.}
create a compiling Lean file with sorry bodies at `PhyXMiniProblems/problem\_phyx\_mini\_0070.lean`. The Lean declarations must preserve the physical quantities, dimensions or dimensional roles, figure labels, governing-law hypotheses, and final relation expressed by this problem.
Use Mathlib/Physlib names found through LeanExplore where available. If a physics API is missing, introduce faithful local abstractions rather than scalar placeholder aliases.

\begin{theorem}[Physics formalization target]
\label{thm:physics:phyx_mini_0070:target}
\lean{PhyXMiniProblems.ProblemPhyXMini0070.problem_phyx_mini_0070}
\uses{def:physics:phyx-mini-0070:phyxminiproblems-problemphyxmini0070-emptytankmetersticksetup, def:physics:phyx-mini-0070:phyxminiproblems-problemphyxmini0070-matchesproblemandprimaryfigure, def:physics:phyx-mini-0070:phyxminiproblems-problemphyxmini0070-hasdepictedphysicalconfiguration, def:physics:phyx-mini-0070:phyxminiproblems-problemphyxmini0070-obeysstraightsightlinegeometry, def:physics:phyx-mini-0070:phyxminiproblems-problemphyxmini0070-roundstowholecentimetermark, def:physics:phyx-mini-0070:phyxminiproblems-problemphyxmini0070-matchesanswerchoice}
The assigned autoformalize agent should translate this physics problem into Lean declarations in the covered file, with theorem and lemma proof bodies written as `by sorry`.
\end{theorem}
\begin{proof}
This is an autoformalization task, not a proof task. Produce statements that can later be proved by the physics prover without weakening the physical model.
\end{proof}

% --- Archon declaration topology begin ---
\section{Declaration topology}
\begin{definition}[Length Quantity]
  \label{def:physics:phyx-mini-0070:phyxminiproblems-problemphyxmini0070-lengthquantity}
  \lean{PhyXMiniProblems.ProblemPhyXMini0070.LengthQuantity}
  A physical length or signed one-dimensional displacement
\end{definition}
\begin{proof}
  This is the declared model definition of Length Quantity.
\end{proof}
\begin{definition}[length Readout]
  \label{def:physics:phyx-mini-0070:phyxminiproblems-problemphyxmini0070-lengthreadout}
  \lean{PhyXMiniProblems.ProblemPhyXMini0070.lengthReadout}
  \uses{def:physics:phyx-mini-0070:phyxminiproblems-problemphyxmini0070-lengthquantity}
  The scalar readout of a dimensionful length in a chosen length unit
\end{definition}
\begin{proof}
  This is the declared model definition of length Readout.
\end{proof}
\begin{definition}[length In Centimeters]
  \label{def:physics:phyx-mini-0070:phyxminiproblems-problemphyxmini0070-lengthincentimeters}
  \lean{PhyXMiniProblems.ProblemPhyXMini0070.lengthInCentimeters}
  \uses{def:physics:phyx-mini-0070:phyxminiproblems-problemphyxmini0070-lengthquantity, def:physics:phyx-mini-0070:phyxminiproblems-problemphyxmini0070-lengthreadout}
  The scalar centimeter readout of a physical length
\end{definition}
\begin{proof}
  This is the declared model definition of length In Centimeters.
\end{proof}
\begin{definition}[Tank Contents]
  \label{def:physics:phyx-mini-0070:phyxminiproblems-problemphyxmini0070-tankcontents}
  \lean{PhyXMiniProblems.ProblemPhyXMini0070.TankContents}
  The material occupying the interior of the tank
\end{definition}
\begin{proof}
  This is the declared model definition of Tank Contents.
\end{proof}
\begin{definition}[Tank Corner]
  \label{def:physics:phyx-mini-0070:phyxminiproblems-problemphyxmini0070-tankcorner}
  \lean{PhyXMiniProblems.ProblemPhyXMini0070.TankCorner}
  Labeled corners of the rectangular tank cross-section
\end{definition}
\begin{proof}
  This is the declared model definition of Tank Corner.
\end{proof}
\begin{definition}[Meter Stick Placement]
  \label{def:physics:phyx-mini-0070:phyxminiproblems-problemphyxmini0070-meterstickplacement}
  \lean{PhyXMiniProblems.ProblemPhyXMini0070.MeterStickPlacement}
  The placement of the meter stick in the primary figure
\end{definition}
\begin{proof}
  This is the declared model definition of Meter Stick Placement.
\end{proof}
\begin{definition}[Angle Reference]
  \label{def:physics:phyx-mini-0070:phyxminiproblems-problemphyxmini0070-anglereference}
  \lean{PhyXMiniProblems.ProblemPhyXMini0070.AngleReference}
  The reference direction used for the displayed sightline angle
\end{definition}
\begin{proof}
  This is the declared model definition of Angle Reference.
\end{proof}
\begin{definition}[Empty Tank Meter Stick Setup]
  \label{def:physics:phyx-mini-0070:phyxminiproblems-problemphyxmini0070-emptytankmetersticksetup}
  \lean{PhyXMiniProblems.ProblemPhyXMini0070.EmptyTankMeterStickSetup}
  \uses{def:physics:phyx-mini-0070:phyxminiproblems-problemphyxmini0070-lengthquantity, def:physics:phyx-mini-0070:phyxminiproblems-problemphyxmini0070-tankcontents, def:physics:phyx-mini-0070:phyxminiproblems-problemphyxmini0070-tankcorner, def:physics:phyx-mini-0070:phyxminiproblems-problemphyxmini0070-meterstickplacement, def:physics:phyx-mini-0070:phyxminiproblems-problemphyxmini0070-anglereference}
  The physical quantities and categorical figure labels for the empty-tank experiment. sightlineBottomRun is the unknown horizontal distance from the stick's zero mark to the point where the straight viewing ray meets the stick; no numerical value for that requested distance is stored in the setup
\end{definition}
\begin{proof}
  This is the declared model definition of Empty Tank Meter Stick Setup.
\end{proof}
\begin{definition}[Matches Problem And Primary Figure]
  \label{def:physics:phyx-mini-0070:phyxminiproblems-problemphyxmini0070-matchesproblemandprimaryfigure}
  \lean{PhyXMiniProblems.ProblemPhyXMini0070.MatchesProblemAndPrimaryFigure}
  \uses{def:physics:phyx-mini-0070:phyxminiproblems-problemphyxmini0070-lengthincentimeters, def:physics:phyx-mini-0070:phyxminiproblems-problemphyxmini0070-emptytankmetersticksetup}
  Problem-statement and primary-figure data. The image reads a 100 cm tank length, a 50 cm tank height, a meter stick extending along the bottom from its zero mark at the left edge, and a sightline grazing the upper-left corner at 30° below the horizontal. The requested bottom mark is not a readout in this predicate
\end{definition}
\begin{proof}
  This is the declared model definition of Matches Problem And Primary Figure.
\end{proof}
\begin{definition}[Has Depicted Physical Configuration]
  \label{def:physics:phyx-mini-0070:phyxminiproblems-problemphyxmini0070-hasdepictedphysicalconfiguration}
  \lean{PhyXMiniProblems.ProblemPhyXMini0070.HasDepictedPhysicalConfiguration}
  \uses{def:physics:phyx-mini-0070:phyxminiproblems-problemphyxmini0070-lengthincentimeters, def:physics:phyx-mini-0070:phyxminiproblems-problemphyxmini0070-emptytankmetersticksetup}
  Positivity, an acute viewing angle, and the condition that the ray's bottom intersection lies on both the tank bottom and the meter stick. These are branch and range conditions, not a numerical assignment of the requested mark
\end{definition}
\begin{proof}
  This is the declared model definition of Has Depicted Physical Configuration.
\end{proof}
\begin{definition}[Obeys Straight Sightline Geometry]
  \label{def:physics:phyx-mini-0070:phyxminiproblems-problemphyxmini0070-obeysstraightsightlinegeometry}
  \lean{PhyXMiniProblems.ProblemPhyXMini0070.ObeysStraightSightlineGeometry}
  \uses{def:physics:phyx-mini-0070:phyxminiproblems-problemphyxmini0070-lengthreadout, def:physics:phyx-mini-0070:phyxminiproblems-problemphyxmini0070-emptytankmetersticksetup}
  Rectilinear propagation through the homogeneous empty tank, expressed as the right-triangle slope relation for the sightline from the upper-left rim to its bottom intersection. It is required in every length unit, so the relation is unit-covariant and is not specialized to the requested centimeter answer
\end{definition}
\begin{proof}
  This is the declared model definition of Obeys Straight Sightline Geometry.
\end{proof}
\begin{definition}[Answer Choice]
  \label{def:physics:phyx-mini-0070:phyxminiproblems-problemphyxmini0070-answerchoice}
  \lean{PhyXMiniProblems.ProblemPhyXMini0070.AnswerChoice}
  Labels of the four displayed multiple-choice answers
\end{definition}
\begin{proof}
  This is the declared model definition of Answer Choice.
\end{proof}
\begin{definition}[answer Mark In Centimeters]
  \label{def:physics:phyx-mini-0070:phyxminiproblems-problemphyxmini0070-answermarkincentimeters}
  \lean{PhyXMiniProblems.ProblemPhyXMini0070.answerMarkInCentimeters}
  \uses{def:physics:phyx-mini-0070:phyxminiproblems-problemphyxmini0070-answerchoice}
  The whole-centimeter meter-stick mark printed beside each answer choice
\end{definition}
\begin{proof}
  This is the declared model definition of answer Mark In Centimeters.
\end{proof}
\begin{definition}[Rounds To Whole Centimeter Mark]
  \label{def:physics:phyx-mini-0070:phyxminiproblems-problemphyxmini0070-roundstowholecentimetermark}
  \lean{PhyXMiniProblems.ProblemPhyXMini0070.RoundsToWholeCentimeterMark}
  \uses{def:physics:phyx-mini-0070:phyxminiproblems-problemphyxmini0070-lengthquantity, def:physics:phyx-mini-0070:phyxminiproblems-problemphyxmini0070-lengthincentimeters}
  The ray intersection rounds to a specified whole-centimeter stick mark
\end{definition}
\begin{proof}
  This is the declared model definition of Rounds To Whole Centimeter Mark.
\end{proof}
\begin{definition}[Matches Answer Choice]
  \label{def:physics:phyx-mini-0070:phyxminiproblems-problemphyxmini0070-matchesanswerchoice}
  \lean{PhyXMiniProblems.ProblemPhyXMini0070.MatchesAnswerChoice}
  \uses{def:physics:phyx-mini-0070:phyxminiproblems-problemphyxmini0070-emptytankmetersticksetup, def:physics:phyx-mini-0070:phyxminiproblems-problemphyxmini0070-answerchoice, def:physics:phyx-mini-0070:phyxminiproblems-problemphyxmini0070-answermarkincentimeters, def:physics:phyx-mini-0070:phyxminiproblems-problemphyxmini0070-roundstowholecentimetermark}
  Agreement of the observed bottom intersection with a displayed choice
\end{definition}
\begin{proof}
  This is the declared model definition of Matches Answer Choice.
\end{proof}
\begin{lemma}[sightline Bottom Run centimeter relation]
  \label{lem:physics:phyx-mini-0070:phyxminiproblems-problemphyxmini0070-sightlinebottomrun-centimeter-relation}
  \lean{PhyXMiniProblems.ProblemPhyXMini0070.sightlineBottomRun_centimeter_relation}
  \uses{def:physics:phyx-mini-0070:phyxminiproblems-problemphyxmini0070-lengthincentimeters, def:physics:phyx-mini-0070:phyxminiproblems-problemphyxmini0070-emptytankmetersticksetup, def:physics:phyx-mini-0070:phyxminiproblems-problemphyxmini0070-matchesproblemandprimaryfigure, def:physics:phyx-mini-0070:phyxminiproblems-problemphyxmini0070-obeysstraightsightlinegeometry}
  The figure readouts and the unit-covariant straight-ray law yield the exact centimeter slope relation before any rounding to a meter-stick mark
\end{lemma}
\begin{proof}
  The conclusion follows from the governing relations and figure readouts named in the statement of sightline Bottom Run centimeter relation.
\end{proof}
% --- Archon declaration topology end ---
% --- Archon physics formalization source end ---
